\section{Method}
Canonical correlation analysis originates in Hotelling (1936) and the two equations that govern the analysis are the following:
\begin{equation} \label{eq:cca_final_b}
    (\Sigma_{21} \Sigma_{11}^{-1} \Sigma_{12} - \rho^2 \Sigma_{22})\mathbf{b} = 0
\end{equation}
\begin{equation} \label{eq:cca_final_a}
    (\Sigma_{12} \Sigma_{22}^{-1} \Sigma_{21} - \rho^2 \Sigma_{11})\mathbf{a} = 0
\end{equation}
where $\Sigma_{21} = \Sigma_{12}^T$. Both equations look similar and have, in fact, the same eigenvalues. And, given the eigenvectors for one of these equations, we can deduce the eigenvectors for the other as will be shown in the next section.

\subsection{Derivation of the canonical correlation analysis equations}
In canonical correlation analysis, we want to maximize correlations between objects that are represented with two data sets. Let us consider two groups of random variables:
\begin{itemize}
    \item Group 1: $\underset{1 \times p}{\mathbf{X}}^{(1)} = (X_1^{(1)}, X_2^{(1)}, \dots, X_p^{(1)})$
    \item Group 2: $\underset{1 \times q}{\mathbf{X}}^{(2)} = (X_1^{(2)}, X_2^{(2)}, \dots, X_q^{(2)})$
\end{itemize}
Sometimes the data in $\mathbf{X}^{(2)}$ and $\mathbf{X}^{(1)}$ are called the dependent and the independent data, respectively. The maximum number of correlations that we can find is then equal to the minimum of the dimensions $p$ and $q$.

Let the directions of optimal correlations for the $\mathbf{X}^{(1)}$ and $\mathbf{X}^{(2)}$ data sets be given by the coefficient vectors $\mathbf{a}$ and $\mathbf{b}$, respectively. We consider the linear combinations $U$ and $V$, defined as follows:
\begin{equation} \label{eq:U_def}
    \underset{1 \times 1}U =  \underset{1 \times p}{\mathbf{X}}^{(1)}\underset{p \times 1} {\mathbf{a}}
\end{equation}
\begin{equation} \label{eq:V_def}
    \underset{1 \times 1}V =  \underset{1 \times q}{\mathbf{X}}^{(2)}\underset{q \times 1} {\mathbf{b}}
\end{equation}

The variables $U$ and $V$ are called the canonical variates. The correlation between them is defined as

\begin{equation}
\rho = \frac{\text{Cov}(U,V)}{\sqrt{\text{Var}(U)\text{Var}(V)}} .
\end{equation}

Since $U = X^{(1)}a$ and $V = X^{(2)}b$, we can express the covariance as

\begin{equation}
\text{Cov}(U,V) 
= \text{Cov}( \mathbf{X}^{(1)} \mathbf{a},  \mathbf{X}^{(2)} \mathbf{b}) 
= \underset{1 \times p}{\mathbf{a}^T} \, \underset{p \times q}{\text{Cov}(\mathbf{X}^{(1)}, \mathbf{X}^{(2)})} \, \underset{q \times 1}{\mathbf{b}}
= \underset{1 \times p}{\mathbf{a}^T} \, \underset{p \times q}{\Sigma_{12}} \, \underset{q \times 1}{\mathbf{b}}
\end{equation}
Similarly, the variances of $U$ and $V$ are

\begin{equation}
\text{Var}(U) = \text{Var}(X^{(1)}a)
= a^T \Sigma_{11} a
\end{equation}

\begin{equation}
\text{Var}(V) = \text{Var}(X^{(2)}b)
= b^T \Sigma_{22} b .
\end{equation}

Substituting these expressions into the definition of correlation yields

\begin{equation} \label{eq:rho_formula}
\rho =
\frac{a^T \Sigma_{12} b}
{\sqrt{a^T \Sigma_{11} a \; b^T \Sigma_{22} b}} .
\end{equation}
Our problem is now finding the directions $\mathbf{b}$ and $\mathbf{a}$ that maximize equation \eqref{eq:rho_formula} above.

We first note that $\rho$ is not affected by a rescaling of $V$ or $U$. Since the choice of rescaling is arbitrary, we therefore maximize equation \eqref{eq:rho_formula} subject to the constraints:
\begin{equation} \label{eq:const_u}
    \text{Var}(U) = \underset{1 \times p}{\mathbf{a}^T} \underset{p \times p}{\Sigma_{11}} \underset{p \times 1}{\mathbf{a}} = 1
\end{equation}
\begin{equation} \label{eq:const_v}
    \text{Var}(V) = \underset{1 \times q}{\mathbf{b}^T} \underset{q \times q}{\Sigma_{22}} \underset{q \times 1}{\mathbf{b}} = 1
\end{equation}
We have made the substitutions where the $\Sigma$'s are covariance matrices. We write the maximization problem in Lagrangian form:
\begin{equation} \label{eq:lagrangian}
    L(\rho_a, \rho_b, \mathbf{a}, \mathbf{b}) = \mathbf{b}^T \Sigma_{21} \mathbf{a} - \frac{\rho_a}{2}(\mathbf{a}^T \Sigma_{11} \mathbf{a} - 1) - \frac{\rho_b}{2}(\mathbf{b}^T \Sigma_{22} \mathbf{b} - 1)
\end{equation}
To solve equation \eqref{eq:lagrangian}, we take partial derivatives with respect to $\mathbf{a}$ and $\mathbf{b}$ and set them to zero:

\begin{equation} \label{eq:deriv_a}
    \frac{\partial L}{\partial \mathbf{a}} = (\mathbf{b}^T\Sigma_{21})^T - \frac{\rho_a}{2} (2 \Sigma_{11} \mathbf{a}) = \Sigma_{12} \mathbf{b} - \rho_a \Sigma_{11} \mathbf{a} = 0 \implies \underset{p \times q}{\Sigma_{12}} \underset{q \times 1}{\mathbf{b}} = \rho_a \underset{p \times p}{\Sigma_{11}} \underset{p \times 1}{\mathbf{a}}
\end{equation}
\begin{equation} \label{eq:deriv_b}
    \frac{\partial L}{\partial \mathbf{b}} = \Sigma_{12}^T \mathbf{a} - \frac{\rho_b}{2} (2 \Sigma_{22} \mathbf{b}) = \Sigma_{21} \mathbf{a} - \rho_b \Sigma_{22} \mathbf{b} = 0 \implies \underset{q \times p}{\Sigma_{21}} \underset{p \times 1}{\mathbf{a}} - \rho_b \underset{q \times q}{\Sigma_{22}} \underset{q \times 1}{\mathbf{b}}
\end{equation}




To identify the relationship between the multipliers $\rho_a$ and $\rho_b$, we left-multiply \eqref{eq:deriv_a} by $\mathbf{a}^T$ and \eqref{eq:deriv_b} by $\mathbf{b}^T$:
\begin{itemize}
    \item From \eqref{eq:deriv_a}: $\mathbf{a}^T \Sigma_{12} \mathbf{b} = \rho_a (\mathbf{a}^T \Sigma_{11} \mathbf{a})$. Given the constraint \eqref{eq:const_u}, we have $\rho_a = \mathbf{a}^T \Sigma_{12} \mathbf{b}$.
    \item From \eqref{eq:deriv_b}: $\mathbf{b}^T \Sigma_{21} \mathbf{a} = \rho_b (\mathbf{b}^T \Sigma_{22} \mathbf{b})$. Given the constraint \eqref{eq:const_v}, we have $\rho_b = \mathbf{b}^T \Sigma_{21} \mathbf{a}$.
\end{itemize}
Since $\mathbf{a}^T \Sigma_{12} \mathbf{b}$ is a scalar, it is equal to its transpose: $(\mathbf{a}^T \Sigma_{12} \mathbf{b})^T = \mathbf{b}^T \Sigma_{12}^T \mathbf{a} = \mathbf{b}^T \Sigma_{21} \mathbf{a}$. This implies that:
\begin{equation} \label{eq:rho_equality}
    \rho_a = \rho_b = \rho = \text{Cov}(U, V)
\end{equation}
Thus, the Lagrange multipliers are identical and equal to the canonical correlation we seek to maximize. Assuming $\Sigma_{11}$ is non-singular, we can rearrange \eqref{eq:deriv_a} to express $\mathbf{a}$ in terms of $\mathbf{b}$:
\begin{equation} \label{eq:a_sol}
\Sigma_{12} \mathbf{b} = \rho_a \Sigma_{11} \mathbf{a} \implies \mathbf{a} = \frac{\Sigma_{11}^{-1} \Sigma_{12} \mathbf{b}}{\rho}
\end{equation}
Substituting \eqref{eq:a_sol} into \eqref{eq:deriv_b} yields the generalized eigenvalue problem for $\mathbf{b}$:
\begin{equation} \label{eq:gen_eig_b}
     \Sigma_{21} \mathbf{a} = \rho_b \Sigma_{22} \mathbf{b} \implies \Sigma_{21} \left( \frac{\Sigma_{11}^{-1} \Sigma_{12} \mathbf{b}}{\rho} \right) = \rho \Sigma_{22} \mathbf{b} \implies (\Sigma_{21} \Sigma_{11}^{-1} \Sigma_{12} - \rho^2 \Sigma_{22})\mathbf{b} = 0
\end{equation}

In an analogous way, we can get the equation for the vectors $\mathbf{a}$ as:
\begin{equation} \label{eq:gen_eig_a}
    (\Sigma_{12} \Sigma_{22}^{-1} \Sigma_{21} - \rho^2 \Sigma_{11})\mathbf{a} = 0
\end{equation}
These equations are so called generalized eigenvalue problems. Special software is needed to solve these equations in a numerically stable and robust manner.

\subsection{Solution of the population CCA problem}
Equation \eqref{eq:gen_eig_b} has the form of a generalized eigenvalue problem
\[
A\mathbf{b} = \lambda B\mathbf{b},
\]
where
\[
A = \Sigma_{21}\Sigma_{11}^{-1}\Sigma_{12}, \quad
B = \Sigma_{22}, \quad
\lambda = \rho^2.
\]
We can consider two cases here: the simple case when we only have the covariance matrices, or, the somewhat more involved case, when we have the original data matrices at our disposal.
If $\Sigma_{22}$ is non-singular, the generalized eigenvalue problem
can be converted into a standard eigenvalue problem by multiplying
\eqref{eq:gen_eig_b} from the left by $\Sigma_{22}^{-1}$, yielding

\begin{equation} \label{eq:eig_standard_b}
    (\Sigma_{22}^{-1} \Sigma_{21} \Sigma_{11}^{-1} \Sigma_{12} - \rho^2 I)\mathbf{b} = 0
\end{equation}

Equation \eqref{eq:eig_standard_b} is a standard eigenvalue problem of the form

\[
M \mathbf{b} = \rho^2 \mathbf{b},
\]

where

\[
M = \Sigma_{22}^{-1} \Sigma_{21} \Sigma_{11}^{-1} \Sigma_{12}.
\]

Thus, the canonical correlations can be obtained by solving the
eigenvalue problem for the matrix $M$. Let

\[
M \mathbf{b}_i = \lambda_i \mathbf{b}_i ,
\]

where $\lambda_i$ are the eigenvalues and $\mathbf{b}_i$ the corresponding
eigenvectors. The canonical correlations are then given by

\[
\rho_i = \sqrt{\lambda_i}.
\]

Once the vectors $\mathbf{b}_i$ are obtained, the corresponding vectors
$\mathbf{a}_i$ can be computed from equation \eqref{eq:deriv_a} as

\[
\mathbf{a}_i =
\frac{\Sigma_{11}^{-1} \Sigma_{12} \mathbf{b}_i}{\rho_i}.
\]

The vectors $\mathbf{a}_i$ and $\mathbf{b}_i$ define the canonical
directions, and the corresponding canonical variates are

\[
U_i = \mathbf{a}_i^T \mathbf{X}^{(1)}, \qquad
V_i = \mathbf{b}_i^T \mathbf{X}^{(2)} .
\]

The number of non–zero canonical correlations is at most
$\min(p,q)$.


\subsection{Solution of the sample CCA problem }
In practical applications, the true population parameters are unknown. We estimate the canonical correlations and weights using the full Linnerud dataset n=20.
\subsubsection{Data Centering}
 Let $\mathbf{X}^{(1)} \in \mathbb{R}^{n \times p}$ and $\mathbf{X}^{(2)} \in \mathbb{R}^{n \times q}$ be the observed data matrices. For CCA to correctly identify relationships based on variance, the data must first be centered:
\begin{equation}
\bar{\mathbf{X}}^{(1)} =
\mathbf{X}^{(1)} -
\begin{pmatrix}
1\\
1\\
\vdots\\
1
\end{pmatrix}
\boldsymbol{\mu}^{(1)},
\quad
\bar{\mathbf{X}}^{(2)} =
\mathbf{X}^{(2)} -
\begin{pmatrix}
1\\
1\\
\vdots\\
1
\end{pmatrix}
\boldsymbol{\mu}^{(2)}
\end{equation}
 where $\boldsymbol{\mu}$ represents the vector of column means. This ensures that the variables have a mean of zero, allowing the subsequent matrix products to represent covariances.
 \subsubsection{Sample Covariance Estimation}
 The sample covariance matrices are estimated from the centered data as follows:\begin{equation}\hat{\Sigma}_{11} = \frac{1}{n-1}\bar{\mathbf{X}}^{(1)T}\bar{\mathbf{X}}^{(1)}, \quad \hat{\Sigma}_{22} = \frac{1}{n-1}\bar{\mathbf{X}}^{(2)T}\bar{\mathbf{X}}^{(2)}, \quad \hat{\Sigma}_{12} = \frac{1}{n-1}\bar{\mathbf{X}}^{(1)T}\bar{\mathbf{X}}^{(2)}
\end{equation}
where $\hat{\Sigma}_{11} \in \mathbb{R}^{p \times p}$ and $\hat{\Sigma}_{22} \in \mathbb{R}^{q \times q}$ represent the internal variance structures, and $\hat{\Sigma}_{12} \in \mathbb{R}^{p \times q}$ captures the cross-set associations.\subsubsection{ Solving the Empirical Eigenvalue Problem}
We solve for the sample canonical weights by substituting the estimates into the standard eigenvalue problem:\begin{equation} \label{eq:sample_eig_b}(\hat{\Sigma}_{22}^{-1} \hat{\Sigma}_{21} \hat{\Sigma}_{11}^{-1} \hat{\Sigma}_{12}) \hat{\mathbf{b}}_i = \hat{\rho}_i^2 \hat{\mathbf{b}}i
\end{equation}
The sample canonical correlations $\hat{\rho}_i$ are the square roots of the eigenvalues. The weight vectors $\hat{\mathbf{a}}_i$ for the exercise set are then recovered:
\begin{equation}
\hat{\mathbf{a}}_i = \frac{\hat{\Sigma}_{11}^{-1} \hat{\Sigma}_{12} \hat{\mathbf{b}}_i}{\hat{\rho}_i}\end{equation}
\subsubsection{ Computation of Sample Scores}
The final canonical scores for the 20 subjects are obtained by projecting the centered data onto the estimated canonical directions:\begin{equation}\hat{\mathbf{u}}_i = \bar{\mathbf{X}}^{(1)} \hat{\mathbf{a}}_i, \qquad \hat{\mathbf{v}}_i = \bar{\mathbf{X}}^{(2)} \hat{\mathbf{b}}_i\end{equation}These scores $\hat{\mathbf{u}}_i, \hat{\mathbf{v}}_i \in \mathbb{R}^{n}$ provide a simplified, one-dimensional representation of each subject's position along the identified latent axes of "Physical Effort" and "Physiological Profile".
\subsection{Hypothesis Testing and Statistical Significance}

After calculating the sample canonical correlations $\hat{\rho}_i$, we must determine whether these associations represent true relationships in the population or are merely products of sampling error.
\subsubsection{Overall Test of Independence}
The primary hypothesis tests whether there is any linear relationship between the two sets of variables:\begin{itemize}\item $H_0: \rho_1 = \rho_2 = \dots = \rho_k = 0$ (The two sets are independent)\item $H_1: \text{At least one } \rho_i \neq 0$ (There is at least one significant association)\end{itemize}where $k = \min(p, q)$.
\subsubsection{Wilks' Lambda Statistic}
The test is based on Wilks' Lambda ($\Lambda$), which measures the proportion of variance not explained by the canonical variables. It is calculated using the sample correlations:\begin{equation}\Lambda_1 = \prod_{i=1}^{k} (1 - \hat{\rho}_i^2)\end{equation}A value of $\Lambda$ near $0$ indicates a strong association, while a value near $1$ suggests independence.
\subsubsection{Bartlett's Chi-square Transformation}
For large samples, Bartlett proposed a transformation of $\Lambda$ to a Chi-square distribution to simplify the calculation of $p$-values. The test statistic is given by:\begin{equation}\chi^2 = - \left( n - 1 - \frac{p + q + 1}{2} \right) \ln \Lambda_1\end{equation}This statistic follows a Chi-square distribution with $\mathbf{df = p \times q}$ degrees of freedom.
\subsubsection{Sequential Testing (Dimension Reduction)}
If the overall $H_0$ is rejected, we conclude that the first pair of canonical variables is significant. We then proceed to test if the remaining correlations are significant by "peeling off" the first one:\begin{equation}\Lambda_m = \prod_{i=m}^{k} (1 - \hat{\rho}_i^2)\end{equation}The degrees of freedom for each subsequent step $m$ are adjusted to:\begin{equation}df = (p - m + 1)(q - m + 1)\end{equation}This sequential process allows us to determine the exact number of significant canonical dimensions to interpret.
\subsubsection{Application to Linnerud Dataset }
For the full Linnerud dataset ($p=3, q=2$), we have:\begin{itemize}\item Step 1: Overall test ($m=1$) with $df = 3 \times 2 = 6$.\item Step 2: Testing the remaining correlation ($m=2$) with $df = (3-1)(2-1) = 2$.\end{itemize}If the $p$-value for Step 1 is less than the significance level (e.g., $\alpha = 0.05$), we reject the null hypothesis and proceed to interpret the first canonical dimension.